% OpenStax Calculus Volume 3 -- Section 5.5 Exercises: Triple Integrals in Cylindrical and Spherical Coordinates
% Exercises reproduced from OpenStax, licensed under CC BY 4.0.
% Source: https://openstax.org/books/calculus-volume-3/pages/5-5-triple-integrals-in-cylindrical-and-spherical-coordinates
\documentclass{exercises}
\title{OpenStax Calculus Volume 3}
\subtitle{Section 5.5 Exercises: Triple Integrals in Cylindrical and Spherical Coordinates}
\sourceurl{https://openstax.org/books/calculus-volume-3/pages/5-5-triple-integrals-in-cylindrical-and-spherical-coordinates}
\author{}
\date{}

\begin{document}
\maketitle


In the following exercises, evaluate the triple integrals $\underset{B}{\iiint}f(x,y,z)\,dV$ over the solid $B$.

\begin{enumerate}[start=1]
\item $f(x,y,z)=z$, $B=\{(x,y,z)|x^{2}+y^{2}\le9,x\ge0,y\ge0,0\le z\le1\}$ \\ \textit{[Figure: A quarter section of a cylinder with height 1 and radius 3.]}
\item $f(x,y,z)=xz^{2}$, $B=\{(x,y,z)|x^{2}+y^{2}\le16,x\ge0,y\le0,-1\le z\le1\}$
\item $f(x,y,z)=xy$, $B=\{(x,y,z)|x^{2}+y^{2}\le1,x\ge0,y\ge0,x\le y,-1\le z\le1\}$ \\ \textit{[Figure: A wedge with radius 1, height 1, and angle pi/4.]}
\item $f(x,y,z)=x^{2}+y^{2}$, $B=\{(x,y,z)|x^{2}+y^{2}\le4,x\ge0,x\le y,0\le z\le3\}$
\item $f(x,y,z)=e^{\sqrt{x^{2}+y^{2}}}$, $B=\{(x,y,z)|1\le x^{2}+y^{2}\le4,y\le0,x\le y\sqrt{3},2\le z\le3\}$
\item $f(x,y,z)=\sqrt{x^{2}+y^{2}}$, $B=\{(x,y,z)|1\le x^{2}+y^{2}\le9,y\le0,0\le z\le1\}$
\item
\begin{enumerate}[label=(\alph*)]
  \item Let $B$ be a cylindrical shell with inner radius $a$, outer radius $b$, and height $c$, where $0<a<b$ and $c>0$. Assume that a function $F$ defined on $B$ can be expressed in cylindrical coordinates as $F(x,y,z)=f(r)+h(z)$, where $f$ and $h$ are differentiable functions. If $\int_{a}^{b}\tilde{f}(r)\,dr=0$ and $\tilde{h}(0)=0$, where $\tilde{f}$ and $\tilde{h}$ are antiderivatives of $f$ and $h$, respectively, show that\\ $\underset{B}{\iiint}F(x,y,z)\,dV=2\pi c(b\tilde{f}(b)-a\tilde{f}(a))+\pi(b^{2}-a^{2})\tilde{h}(c)$.
  \item Use the previous result to show that $\underset{B}{\iiint}(z+\sin \sqrt{x^{2}+y^{2}})\,dx\,dy\,dz=6\pi^{2}(\pi-2)$, where $B$ is a cylindrical shell with inner radius $\pi$, outer radius $2\pi$, and height $2$.
\end{enumerate}
\item
\begin{enumerate}[label=(\alph*)]
  \item Let $B$ be a cylindrical shell with inner radius $a$, outer radius $b$, and height $c$, where $0<a<b$ and $c>0$. Assume that a function $F$ defined on $B$ can be expressed in cylindrical coordinates as $F(x,y,z)=f(r)g(\theta)h(z)$, where $f$, $g$, and $h$ are differentiable functions. If $\int_{a}^{b}\tilde{f}(r)\,dr=0$, where $\tilde{f}$ is an antiderivative of $f$, show that\\ $\underset{B}{\iiint}F(x,y,z)\,dV=[b\tilde{f}(b)-a\tilde{f}(a)]\,[\tilde{g}(2\pi)-\tilde{g}(0)]\,[\tilde{h}(c)-\tilde{h}(0)]$,\\ where $\tilde{g}$ and $\tilde{h}$ are antiderivatives of $g$ and $h$, respectively.
  \item Use the previous result to show that $\underset{B}{\iiint}z\sin \sqrt{x^{2}+y^{2}}\,dx\,dy\,dz=-12\pi^{2}$, where $B$ is a cylindrical shell with inner radius $\pi$, outer radius $2\pi$, and height $2$.
\end{enumerate}
\end{enumerate}

In the following exercises, the boundaries of the solid $E$ are given in cylindrical coordinates. Let $f\,(r,\theta,z)$ be the corresponding function in cylindrical coordinates.


(a) Define the region in cylindrical coordinates.


(b) Convert the integral $\underset{E}{\iiint}g(x,y,z)\,dV$ to cylindrical coordinates.

\begin{enumerate}[resume]
\item $E$ is inside the right circular cylinder $r=4\sin \theta$, above the $r\theta$-plane, and inside the sphere $r^{2}+z^{2}=16$.
\item $E$ is inside the right circular cylinder $r=\cos \theta$, above the $r\theta$-plane, and inside the sphere $r^{2}+z^{2}=9$.
\item $E$ is located in the first octant and is bounded by the circular paraboloid $z=9-3r^{2}$, the cylinder $r=\sqrt{3}$, and the plane $r(\cos \theta+\sin \theta)=20-z$.
\item $E$ is located in the first octant outside the circular paraboloid $z=10-2r^{2}$ and inside the cylinder $r=\sqrt{5}$ and is bounded also by the planes $z=20$ and $\theta=\frac{\pi}{4}$.
\end{enumerate}

In the following exercises, the function $g$ and region $E$ are given in rectangular coordinates.


(a) Express the region $E$ and the function $g$ in cylindrical coordinates. Let $f(r,\theta,z)$ be the corresponding function in cylindrical coordinates.


(b) Convert the integral $\underset{E}{\iiint}g(x,y,z)\,dV$ to cylindrical coordinates and evaluate it.

\begin{enumerate}[resume]
\item $g(x,y,z)=\frac{1}{x+3}$, $E=\{(x,y,z)|0\le x^{2}+y^{2}\le9,x\ge0,y\ge0,0\le z\le x+3\}$
\item $g(x,y,z)=x^{2}+y^{2}$, $E=\{(x,y,z)|0\le x^{2}+y^{2}\le4,y\ge0,0\le z\le3-x\}$
\item $g(x,y,z)=x$, $E=\{(x,y,z)|1\le y^{2}+z^{2}\le9,0\le x\le9-y^{2}-z^{2}\}$
\item $g(x,y,z)=y$, $E=\{(x,y,z)|1\le x^{2}+z^{2}\le9,0\le y\le9-x^{2}-z^{2}\}$
\end{enumerate}

In the following exercises, find the volume of the solid $E$ whose boundaries are given in rectangular coordinates.

\begin{enumerate}[resume]
\item $E$ is above the $xy$-plane, inside the cylinder $x^{2}+y^{2}=1$, and below the plane $z=1$.
\item $E$ is below the plane $z=1$ and inside the paraboloid $z=x^{2}+y^{2}$.
\item $E$ is bounded by the circular cone $z=\sqrt{x^{2}+y^{2}}$ and $z=1$.
\item $E$ is located above the $xy$-plane, below $z=1$, outside the one-sheeted hyperboloid $x^{2}+y^{2}-z^{2}=1$, and inside the cylinder $x^{2}+y^{2}=2$.
\item $E$ is located inside the cylinder $x^{2}+y^{2}=1$ and between the circular paraboloids $z=1-x^{2}-y^{2}$ and $z=x^{2}+y^{2}$.
\item $E$ is located inside the sphere $x^{2}+y^{2}+z^{2}=1$, above the $xy$-plane, and inside the circular cone $z=\sqrt{x^{2}+y^{2}}$.
\item $E$ is located inside the circular cone $x^{2}+y^{2}={(z-1)}^{2}$ and between the planes $z=0$ and $z=2$.
\item $E$ is located inside the cylinder $x^{2}+y^{2}=1$, above the circular cone $z=1-\sqrt{x^{2}+y^{2}}$, and below the circular paraboloid $z=1+x^{2}+y^{2}$, and between the planes $z=0$ and $z=2$.
\item \techrequired Use a computer algebra system (CAS) to graph the solid whose volume is given by the iterated integral in cylindrical coordinates $\int_{-\pi/2}^{\pi/2}\,\int_{0}^{1}\,\int_{r^{2}}^{r}r\,dz\,dr\,d\theta$. Find the volume $V$ of the solid. Round your answer to four decimal places.
\item \techrequired Use a CAS to graph the solid whose volume is given by the iterated integral in cylindrical coordinates $\int_{0}^{\pi/2}\,\int_{0}^{1}\,\int_{r^{4}}^{r}r\,dz\,dr\,d\theta$. Find the volume $V$ of the solid. Round your answer to four decimal places.
\item Convert the integral $\int_{0}^{1}\,\int_{-\sqrt{1-y^{2}}}^{\sqrt{1-y^{2}}}\,\int_{x^{2}+y^{2}}^{\sqrt{x^{2}+y^{2}}}xz\,dz\,dx\,dy$ into an integral in cylindrical coordinates.
\item Convert the integral $\int_{0}^{2}\,\int_{0}^{y}\,\int_{0}^{1}(xy+z)\,dz\,dx\,dy$ into an integral in cylindrical coordinates.
\end{enumerate}

In the following exercises, evaluate the triple integral $\underset{B}{\iiint}f(x,y,z)\,dV$ over the solid $B$.

\begin{enumerate}[resume]
\item $f(x,y,z)=1$, $B=\{(x,y,z)|x^{2}+y^{2}+z^{2}\le90,z\ge0\}$ \\ \textit{[Figure: A filled-in half-sphere with radius 3 times the square root of 10.]}
\item $f(x,y,z)=1-\sqrt{x^{2}+y^{2}+z^{2}}$, $B=\{(x,y,z)|x^{2}+y^{2}+z^{2}\le9,y\ge0,z\ge0\}$ \\ \textit{[Figure: A quarter section of an ovoid with height 8, width 8 and length 18.]}
\item $f(x,y,z)=\sqrt{x^{2}+y^{2}}$, $B$ is bounded above by the half-sphere $x^{2}+y^{2}+z^{2}=9$ with $z\ge0$ and below by the cone $z^{2}=x^{2}+y^{2}$.
\item $f(x,y,z)=z$, $B$ is bounded above by the half-sphere $x^{2}+y^{2}+z^{2}=16$ with $z\ge0$ and below by the cone $z^{2}=x^{2}+y^{2}$.
\item Show that if $F(\rho,\theta,\phi)=f(\rho)g(\theta)h(\phi)$ is a continuous function on the spherical box $B=\{(\rho,\theta,\phi)|a\le \rho \le b,\alpha \le \theta \le \beta,\gamma \le \phi \le \psi \}$, then \[\underset{B}{\iiint}F\,dV=(\int_{a}^{b}\rho^{2}f(\rho)\,d\rho)(\int_{\alpha}^{\beta}g(\theta)\,d\theta)(\int_{\gamma}^{\psi}h(\phi)\sin \phi \,d\phi).\]
\item
\begin{enumerate}[label=(\alph*)]
  \item A function $F$ is said to have spherical symmetry if it depends on the distance to the origin only, that is, it can be expressed in spherical coordinates as $F(x,y,z)=f(\rho)$, where $\rho=\sqrt{x^{2}+y^{2}+z^{2}}$. Show that\\ $\underset{B}{\iiint}F(x,y,z)\,dV=2\pi \int_{a}^{b}\rho^{2}f(\rho)\,d\rho$,\\ where $B$ is the region between the upper concentric hemispheres of radii $a$ and $b$ centered at the origin, with $0<a<b$ and $F$ a spherical function defined on $B$.
  \item Use the previous result to show that $\underset{B}{\iiint}(x^{2}+y^{2}+z^{2})\sqrt{x^{2}+y^{2}+z^{2}}\,dV=21\pi$, where\\ $B=\{(x,y,z)|1\le x^{2}+y^{2}+z^{2}\le4,z\ge0\}$.
\end{enumerate}
\item
\begin{enumerate}[label=(\alph*)]
  \item Let $B$ be the region between the upper concentric hemispheres of radii $a$ and $b$ centered at the origin and situated in the first octant, where $0<a<b$. Consider F a function defined on B whose form in spherical coordinates $(\rho,\theta,\phi)$ is $F(x,y,z)=f(\rho)\cos \phi$. Show that if $g(a)=g(b)=0$ and $\int_{a}^{b}h(\rho)\,d\rho=0$, then\\ $\underset{B}{\iiint}F(x,y,z)\,dV=\frac{\pi}{2}[ah(a)-bh(b)]$,\\ where $g$ is an antiderivative of $f$ and $h$ is an antiderivative of $g$.
  \item Use the previous result to show that $\underset{B}{\iiint}\frac{z\cos \sqrt{x^{2}+y^{2}+z^{2}}}{\sqrt{x^{2}+y^{2}+z^{2}}}\,dV=\frac{3\pi^{2}}{2}$, where $B$ is the region between the upper concentric hemispheres of radii $\pi$ and $2\pi$ centered at the origin and situated in the first octant.
\end{enumerate}
\end{enumerate}

In the following exercises, the function $g$ and region $E$ are given in rectangular coordinates.


(a) Express the region $E$ and the function $g$ in spherical coordinates. Let $f(\rho,\theta,\phi)$ be the corresponding function in spherical coordinates.


(b) Convert the integral $\underset{E}{\iiint}g(x,y,z)\,dV$ to spherical coordinates and evaluate it.

\begin{enumerate}[resume]
\item $g(x,y,z)=z$; $E=\{(x,y,z)|0\le x^{2}+y^{2}+z^{2}\le1,z\ge0\}$
\item $g(x,y,z)=x+y$; $E=\{(x,y,z)|1\le x^{2}+y^{2}+z^{2}\le4,z\ge0,y\ge0\}$
\item $g(x,y,z)=2xy$; $E=\{(x,y,z)|\sqrt{x^{2}+y^{2}}\le z\le \sqrt{1-x^{2}-y^{2}},x\ge0,y\ge0\}$
\item $g(x,y,z)=z$; $E=\{(x,y,z)|x^{2}+y^{2}+z^{2}-2z\le0,\sqrt{x^{2}+y^{2}}\le z\}$
\end{enumerate}

In the following exercises, find the volume of the solid $E$ whose boundaries are given in rectangular coordinates.

\begin{enumerate}[resume]
\item $E=\{(x,y,z)|\sqrt{x^{2}+y^{2}}\le z\le \sqrt{16-x^{2}-y^{2}},x\ge0,y\ge0\}$
\item $E=\{(x,y,z)|x^{2}+y^{2}+z^{2}-2z\le0,\sqrt{x^{2}+y^{2}}\le z\}$
\item Use spherical coordinates to find the volume of the solid situated inside the sphere $\rho=1$ and outside the sphere $\rho=\cos \phi$, with $\phi \in[0,\frac{\pi}{2}]$.
\item Use spherical coordinates to find the volume of the ball $\rho \le3$ that is situated between the cones $\phi=\frac{\pi}{4}$ and $\phi=\frac{\pi}{3}$.
\item Convert the integral $\int_{-4}^{4}\,\int_{-\sqrt{16-y^{2}}}^{\sqrt{16-y^{2}}}\,\int_{-\sqrt{16-x^{2}-y^{2}}}^{\sqrt{16-x^{2}-y^{2}}}(x^{2}+y^{2}+z^{2})\,dz\,dx\,dy$ into an integral in spherical coordinates.
\item Convert the integral $\int_{0}^{4}\,\int_{0}^{\sqrt{16-x^{2}}}\,\int_{-\sqrt{16-x^{2}-y^{2}}}^{\sqrt{16-x^{2}-y^{2}}}{(x^{2}+y^{2}+z^{2})}^{2}\,dz\,dy\,dx$ into an integral in spherical coordinates.
\item Convert the integral $\int_{-2\sqrt{2}}^{2\sqrt{2}}\,\int_{-\sqrt{8-x^{2}}}^{\sqrt{8-x^{2}}}\,\int_{\sqrt{x^{2}+y^{2}}}^{\sqrt{16-x^{2}-y^{2}}}\,dz\,dy\,dx$ into an integral in spherical coordinates and evaluate it.
\item \techrequired Use a CAS to graph the solid whose volume is given by the iterated integral in spherical coordinates $\int_{\pi/2}^{\pi}\,\int_{5\pi/6}^{\pi/6}\,\int_{0}^{2}\rho^{2}\sin \phi \,d\rho \,d\phi \,d\theta$. Find the volume $V$ of the solid. Round your answer to three decimal places.
\item \techrequired Use a CAS to graph the solid whose volume is given by the iterated integral in spherical coordinates as $\int_{0}^{2\pi}\,\int_{\pi/4}^{3\pi/4}\,\int_{0}^{1}\rho^{2}\sin \phi \,d\rho \,d\phi \,d\theta$. Find the volume $V$ of the solid. Round your answer to three decimal places.
\item \techrequired Use a CAS to evaluate the integral $\underset{E}{\iiint}(x^{2}+y^{2})\,dV$ where $E$ lies above the paraboloid $z=x^{2}+y^{2}$ and below the plane $z=3y$.
\item \techrequired
\begin{enumerate}[label=(\alph*)]
  \item Evaluate the integral $\underset{E}{\iiint}e^{\sqrt{x^{2}+y^{2}+z^{2}}}\,dV$, where $E$ is bounded by the spheres $4x^{2}+4y^{2}+4z^{2}=1$ and $x^{2}+y^{2}+z^{2}=1$.
  \item Use a CAS to find an approximation of the previous integral. Round your answer to two decimal places.
\end{enumerate}
\item Express the volume of the solid inside the sphere $x^{2}+y^{2}+z^{2}=16$ and outside the cylinder $x^{2}+y^{2}=4$ as triple integrals in cylindrical coordinates and spherical coordinates.
\item Express the volume of the solid inside the sphere $x^{2}+y^{2}+z^{2}=16$ and outside the cylinder $x^{2}+y^{2}=4$ that is located in the first octant as triple integrals in cylindrical coordinates and spherical coordinates.
\item The power emitted by an antenna has a power density per unit volume given in spherical coordinates by $p(\rho,\theta,\phi)=\frac{P_{0}}{\rho^{2}}\cos^{2}\theta \sin^{4}\phi$, where $P_{0}$ is a constant with units in watts. The total power within a sphere $B$ of radius $r$ meters is defined as $P=\underset{B}{\iiint}p(\rho,\theta,\phi)\,dV$. Find the total power $P$ within a sphere of radius 20 meters.
\item Use the preceding exercise to find the total power within a sphere $B$ of radius 5 meters when the power density per unit volume is given by $p(\rho,\theta,\phi)=\frac{30}{\rho^{2}}\cos^{2}\theta \sin^{4}\phi$.
\item A charge cloud contained in a sphere $B$ of radius $r$ centimeters centered at the origin has its charge density given by $q(x,y,z)=k\sqrt{x^{2}+y^{2}+z^{2}}\frac{\mu \,\text{C}}{\text{cm}^{3}}$, where $k>0$. The total charge contained in $B$ is given by $Q=\underset{B}{\iiint}q(x,y,z)\,dV$. Find the total charge $Q$.
\item Use the preceding exercise to find the total charge cloud contained in the unit sphere if the charge density is $q(x,y,z)=20\sqrt{x^{2}+y^{2}+z^{2}}\frac{\mu \,\text{C}}{\text{cm}^{3}}$.
\end{enumerate}

\end{document}
